% Ch4 WS1 -- solutions, electrolytes, molarity, dilution. Block 1.
\documentclass{apchem-worksheet}

\apcunitword{Chapter}
\apcunit{4}
\apcunittitle{Types of Chemical Reactions and Solution Stoichiometry}
\apcdoctitle{Worksheet 1 \textbullet\ Solutions and Electrolytes}
\apcsubtitle{Zumdahl \S4.1--4.3 \textbullet\ strong is not the same word as concentrated}

\begin{document}
\wsheader[Show the setup for every calculation. For classification
questions, state what happens to the solute in water --- not just the
label.]

\problem[]{Classify each as a strong electrolyte, weak electrolyte, or
nonelectrolyte, and state what species are present in solution:}
\begin{parts}
  \item \ce{KBr}: \answerblank[6cm]{strong; \ce{K+} and \ce{Br-}}
  \item \ce{HC2H3O2}: \answerblank[6cm]{weak; mostly intact molecules}
  \item \ce{C6H12O6} (glucose): \answerblank[6cm]{nonelectrolyte; molecules}
  \item \ce{HNO3}: \answerblank[6cm]{strong; \ce{H+} and \ce{NO3-}}
  \item \ce{NH3}: \answerblank[6cm]{weak; mostly \ce{NH3} molecules}
\end{parts}

\problem[]{Two beakers each hold \SI{0.10}{\Molar} solutions. Beaker A
(\ce{HCl}) lights a conductivity bulb brightly; beaker B (\ce{HF}) lights it
dimly.}
\begin{parts}
  \item Explain the difference. \answerlines[2]{\ce{HCl} is a strong acid
        and ionizes essentially completely, giving a high ion
        concentration; \ce{HF} is weak and only slightly ionized, so far
        fewer ions carry current.}
  \item A student concludes ``beaker A is more concentrated.'' Evaluate.
        \answerlines[2]{Wrong --- both are 0.10 M. Strong describes the
        fraction that ionizes, not the amount of solute present.}
\end{parts}

\problem[]{Explain at the particle level why water dissolves \ce{NaCl} but
\ce{CCl4} does not. \answerlines[3]{Water is polar, so its oxygen ends
attract \ce{Na+} and its hydrogen ends attract \ce{Cl-}; these ion--dipole
attractions release enough energy to overcome the lattice. \ce{CCl4} is
nonpolar and offers only weak dispersion forces, which cannot pay for
breaking the lattice.}}

\problem[]{Molarity calculations:}
\begin{parts}
  \item Molarity of \SI{29.2}{\gram} \ce{NaCl} in \SI{500.}{\milli\liter}:
        \answerblank[3cm]{\SI{1.00}{\Molar}}
  \item Mass of \ce{KMnO4} ($M = 158.03$) in \SI{250.}{\milli\liter} of
        \SI{0.0200}{\Molar}: \answerblank[3cm]{\SI{0.790}{\gram}}
  \item Moles in \SI{40.0}{\milli\liter} of \SI{0.125}{\Molar}:
        \answerblank[3cm]{\SI{5.00e-3}{\mole}}
  \item Volume of \SI{0.250}{\Molar} containing \SI{0.0500}{\mole}:
        \answerblank[3cm]{\SI{200.}{\milli\liter}}
\end{parts}

\problem[]{Dilutions:}
\begin{parts}
  \item \SI{15.0}{\milli\liter} of \SI{18.0}{\Molar} \ce{H2SO4} diluted to
        \SI{500.}{\milli\liter}: \answerblank[3cm]{\SI{0.540}{\Molar}}
  \item Volume of \SI{12.0}{\Molar} \ce{HCl} needed for
        \SI{100.}{\milli\liter} of \SI{0.500}{\Molar}:
        \answerblank[3cm]{\SI{4.17}{\milli\liter}}
  \item State the safety rule for diluting concentrated acid, and why.
        \answerlines[2]{Add acid to water, never water to acid --- dilution
        is strongly exothermic, and adding water to concentrated acid can
        boil and spatter it.}
\end{parts}

\problem[]{Ion concentrations in strong electrolyte solutions:}
\begin{parts}
  \item $[\ce{K+}]$ and $[\ce{SO4^2-}]$ in \SI{0.30}{\Molar} \ce{K2SO4}:
        \answerblank[5cm]{\SI{0.60}{\Molar}; \SI{0.30}{\Molar}}
  \item $[\ce{NO3-}]$ in \SI{0.15}{\Molar} \ce{Fe(NO3)3}:
        \answerblank[3cm]{\SI{0.45}{\Molar}}
  \item Total ion concentration in \SI{0.20}{\Molar} \ce{Na3PO4}:
        \answerblank[3cm]{\SI{0.80}{\Molar}}
\end{parts}

\problem[]{Describe, step by step, how to prepare \SI{250.0}{\milli\liter}
of \SI{0.100}{\Molar} \ce{K2Cr2O7} ($M = 294.20$) from the solid.
\answerlines[3]{$n = 0.0250$ mol; $m = 0.0250 \times 294.20 =
\SI{7.36}{\gram}$. Weigh it, dissolve in about \SI{150}{\milli\liter} of
water, transfer quantitatively to a \SI{250}{\milli\liter} volumetric flask,
rinse the beaker into the flask, add water to the mark, and invert to
mix.}}

\begin{spiralstrip}{Chapters 2--3 \textbullet\ spiral}
  \item Name \ce{Cu(NO3)2}: \answerblank[4.5cm]{copper(II) nitrate}
  \item Moles in \SI{44.0}{\gram} \ce{CO2}: \answerblank[2.6cm]{\SI{1.00}{\mole}}
  \item Empirical formula of \ce{C6H12}: \answerblank[2.4cm]{\ce{CH2}}
  \item Limiting reactant idea: theoretical yield comes from which
        reactant? \answerblank[3.4cm]{the limiting one}
  \item Electrons in \ce{S^2-}: \answerblank[2cm]{18}
\end{spiralstrip}

\end{document}
